\documentclass[11pt]{article}
\usepackage{amsmath, amssymb, amsthm}
\usepackage{geometry}
\usepackage{hyperref}
\geometry{margin=1in}

\newtheorem{theorem}{Theorem}
\newtheorem{lemma}[theorem]{Lemma}
\newtheorem{proposition}[theorem]{Proposition}
\newtheorem{corollary}[theorem]{Corollary}
\theoremstyle{definition}
\newtheorem{definition}[theorem]{Definition}
\newtheorem{remark}[theorem]{Remark}
\newtheorem{observation}[theorem]{Observation}

\title{Sandvik Monodromic Endoscopy vs.\ Block Decomposition Rule B:\\
A Shape Mismatch at $n=6,\ j=2$}
\author{Clio}
\date{2026-04-21 (Prove session)}

\begin{document}
\maketitle

\begin{abstract}
We test the conjecture, raised in \texttt{PROVE.md} of 2026-04-21, that Block
Decomposition Rule B for the staircase matrix
$M_R^{(k)}[D_n, D_n]$ is the shape-level realisation of Sandvik's monodromic Hecke
endoscopic reduction (arXiv:2604.15582). At the smallest nontrivial test case
$n=6, k=4$ (i.e., the base of the $j=2$ Hirota identity), we show that the Rule B
decomposition is the standard parabolic right-coset decomposition by the parabolic
subgroup $W_k = \langle s_1, \ldots, s_k\rangle \cong S_{k+1}$. Its block sizes
$\{30, 24, 18, 12, 6\}$ are indexed by the value of $w(n)$ and not by a torus
character $\zeta$, so they cannot be the orders of endoscopic Weyl subgroups
$W_\zeta \subseteq S_6$ — in particular there is no Young subgroup of $S_6$ of
orders $30$, $18$, or $12$ at all. The principal block
$(w(6)=1,\text{ size }30)$ does coincide with $M_R^{(4)}[D_5, D_5]$ as a polynomial
matrix, which is the parabolic restriction $S_6 \to S_5$; this agrees with
Sandvik's endoscopic reduction at $\zeta = (z,z,z,z,z,z')$ in the \emph{target} group
$S_5$, but the $D$-restriction kills the $|S_5|/|D_5| = 4$-fold categorical
multiplicity that Sandvik's theorem would supply. The conclusion is that Sandvik's
monodromic framework does not provide a direct categorical lift of Rule B or an
a-priori palindromy of $p_\lambda(q)$ in the staircase Rosetta stone, and the
palindromy should instead be sought in Kazhdan--Lusztig / W-graph positivity on
the parabolic Hecke module $\mathcal{H}_q / \mathcal{H}_q \cdot W_J$ for the
pairing parabolic $W_J = \langle s_1, s_3, s_5\rangle$.
\end{abstract}

\section{Setup}

Throughout, $W = S_n$ with simple generators $s_1, \ldots, s_{n-1}$, and
$\mathcal{H}_q = \mathcal{H}_q(W)$ is its Iwahori--Hecke algebra with
$T_i^2 = (q-1)T_i + q$. We write $W_k = \langle s_1, \ldots, s_k\rangle \cong S_{k+1}$
for the left parabolic on the first $k+1$ positions, and
$W_J = \langle s_1, s_3, s_5, \ldots\rangle \cong (\mathbb{Z}/2)^{\lfloor n/2\rfloor}$
for the pairing parabolic. The descent-paired set is
$$D_n = \left\{w \in S_n : w(2i-1) > w(2i) \text{ for } 1 \le i \le \lfloor n/2\rfloor\right\},$$
a set of minimum-length coset representatives for $S_n / W_J$ with
$|D_n| = n!/2^{\lfloor n/2\rfloor}$.

The partial staircase is
$$\Pi^{(k)} = B_1 B_2 \cdots B_k \in \mathcal{H}_q, \qquad
B_j = (1+T_j)(1+T_{j-1})\cdots(1+T_1),$$
and the right-multiplication matrix $M_R^{(k)}[D_n, D_n] \in \mathrm{Mat}_{|D_n|}(\mathbb{Z}[q])$
has entries $(M_R^{(k)})_{u,v} = $ coefficient of $T_u$ in $T_v \cdot \Pi^{(k)}$, for $u, v \in D_n$.

Block Decomposition Rule B \cite{memory:block-decomposition-rule-b} asserts that at
level $k$ the matrix $M_R^{(k)}[D_n, D_n]$ is exactly block-diagonal under the
partition of $D_n$ by the \emph{suffix} $(w(k+2), \ldots, w(n))$.

\section{Rule B is parabolic}

\begin{lemma}[Rule B $=$ right $W_k$-coset decomposition]\label{lem:rule-b-parabolic}
For $1 \le k \le n-1$, the partition of $S_n$ by the tuple $(w(k+2), \ldots, w(n))$
is the partition by right $W_k$-cosets:
$$v \sim v' \iff v^{-1} v' \in W_k \iff v(i) = v'(i) \text{ for all } i \ge k+2.$$
Consequently, the Rule B block decomposition of $M_R^{(k)}[D_n, D_n]$ is the
intersection of $D_n$ with the right $W_k$-cosets of $S_n$. The partition is invariant
under $M_R^{(k)}$ because $\Pi^{(k)} \in \mathcal{H}(W_k) \subset \mathcal{H}_q$.
\end{lemma}

\begin{proof}
The generators $s_1, \ldots, s_k$ act on one-line notation by swapping adjacent
positions in $\{1, \ldots, k+1\}$ only; every element of $W_k$ is a composition of
such swaps. Hence for $v \in S_n$ and $\sigma \in W_k$ one has $(v\sigma)(i) = v(i)$
for $i \ge k+2$. Conversely, if $v, v'$ agree at positions $k+2, \ldots, n$ then they
have the same values in those positions and differ only in the arrangement of the
remaining values across positions $1, \ldots, k+1$: the rearrangement is
$\sigma = v^{-1} v' \in W_k$. For the matrix statement, $\Pi^{(k)}$ is a polynomial in
$T_1, \ldots, T_k$ only, so right-multiplication by $\Pi^{(k)}$ on $T_v$ produces
linear combinations of $T_{v\sigma}$ for $\sigma \in W_k$ — exactly the right
$W_k$-coset of $v$.
\end{proof}

\begin{proposition}[Block sizes at $n=6, k=4$]\label{prop:block-sizes}
Partitioning $D_6$ by $w(6)$ yields five blocks of sizes
$$(n-c)\cdot|D_{n-2}|\Big|_{n=6} = (6-c)\cdot 6 \in \{30, 24, 18, 12, 6\}
\qquad\text{for}\qquad c = 1, 2, 3, 4, 5.$$
In particular the principal block $c=1$ has size $|D_5| = 30$.
\end{proposition}

\begin{proof}
A permutation $w \in D_6$ with $w(6)=c$ is determined by (i) the value $w(5) \in \{c+1, \ldots, 6\}$ (since $w(5)>w(6)=c$ by the descent-pair condition on $(5,6)$); this gives $6-c$ choices; (ii) a descent-paired arrangement of the remaining four values across positions $1,\ldots,4$, giving $|D_4| = 6$ possibilities. Computational verification at $n=6, k=4$ in SymPy confirmed the exact sizes $30, 24, 18, 12, 6$.
\end{proof}

\begin{proposition}[Principal block identity]\label{prop:principal}
Let $P \subset D_6$ be the principal Rule B block at $n=6, k=4$ (the 30 perms with
$w(6)=1$). Write each $w \in P$ as $w = (w(1),\ldots,w(5), 1)$ and consider the
prefix $w|_{1..5}$ as a permutation of $\{2,\ldots,6\}$, then relabel
$i \mapsto i-1$ to obtain a permutation of $\{1,\ldots,5\}$. The resulting map
$P \to D_5$ is a bijection, and the induced reordering of the submatrix
$M_R^{(4)}[P, P]$ equals $M_R^{(4)}[D_5, D_5]$ entrywise as polynomials in $q$.
\end{proposition}

\begin{proof}
Bijectivity is immediate from the fact that $(w(1),\ldots,w(5))$ ranges over all
descent-paired permutations of the 5-element alphabet $\{2,\ldots,6\}$ as $w$ ranges
over $P$. The matrix equality was verified symbolically: building
$M_R^{(4)}[D_6, D_6]$ as a $90\times 90$ matrix over $\mathbb{Q}(q)$, extracting
the principal sub-block (reordered by the bijection), and subtracting
$M_R^{(4)}[D_5, D_5]$ (a $30 \times 30$ matrix over $\mathbb{Q}(q)$) produced the zero
matrix. The script is \texttt{sandvik-verify.py}.
\end{proof}

\begin{remark}
Proposition \ref{prop:principal} restates the self-similarity noted in \cite{memory:block-decomposition-rule-b}: the Rule B principal block at $(n, k) = (6, 4)$ is literally the $n=5$ version of the same construction. This is parabolic self-similarity, not a monodromic reduction.
\end{remark}

\section{Sandvik's endoscopic reduction}

We recall the relevant structure of Sandvik's paper \cite{sandvik}, which
categorifies the monodromic Hecke algebra of a complex reductive group $G$ with
maximal torus $T$. The monodromic 2-category $\mathcal{H}^\zeta(G)$ at monodromy
$\zeta \in T$ is a Hecke-type category whose Grothendieck group at generic $q$ has
dimension $|W_\zeta|$, where
$$W_\zeta = \{w \in W : w(\zeta) = \zeta\} \subset W$$
is the \emph{endoscopic Weyl subgroup}. Sandvik's main theorem is, schematically,
$$\mathcal{H}^\zeta(G) \;\simeq\; \mathcal{H}^{\text{unip}}(G_\zeta),$$
where $G_\zeta$ is the centralizer of $\zeta$ in $G$ and $\mathcal{H}^{\text{unip}}$
is the unipotent Hecke category.

For $G = GL_n$ with torus $T = (\mathbb{C}^\times)^n$, a character
$\zeta = (\zeta_1, \ldots, \zeta_n)$ has endoscopic subgroup $W_\zeta \cong \prod_i S_{m_i}$
a Young subgroup (where $m_i$ are the multiplicities of distinct torus values);
every reflection subgroup of $S_n$ arises this way.

\begin{lemma}[Young-subgroup orders in $S_6$]\label{lem:young}
The orders of proper Young subgroups $\prod S_{n_i} \subset S_6$ are
$$\{1, 2, 4, 6, 8, 12, 24, 36, 48, 120\} \cup \{720\}.$$
In particular, $30$, $18$ do not appear; $12, 24$ appear but as $|S_3 \times S_2 \times S_1| = 12$,
$|S_4 \times S_1 \times S_1| = 24$ respectively.
\end{lemma}

\begin{proof}
Enumerating compositions of $6$: $6 = 6, 5+1, 4+2, 4+1+1, 3+3, 3+2+1, 3+1+1+1, 2+2+2, 2+2+1+1, 2+1+1+1+1, 1+\cdots+1$. The corresponding orders are
$720, 120, 48, 24, 36, 12, 6, 8, 4, 2, 1$.
\end{proof}

\section{Shape mismatch}

\begin{theorem}\label{thm:mismatch}
At $n=6, k=4, j=2$, the Rule B block decomposition of $M_R^{(4)}[D_6, D_6]$ is
\emph{not} the Sandvik endoscopic decomposition at any monodromy $\zeta \in T$, in the
following sense: the multiset of Rule B block sizes,
$$\{30, 24, 18, 12, 6\},$$
is not a sub-multiset of $\{|W_\zeta| : \zeta \in T\}$ — the Young-subgroup orders in
$S_6$. Specifically, three of the five sizes ($30, 18, 12$) are not Young-subgroup
orders of $S_6$ at all.
\end{theorem}

\begin{proof}
Direct from Lemma \ref{lem:young} and Proposition \ref{prop:block-sizes}.
\end{proof}

\begin{corollary}[Against the $D$-restricted match]\label{cor:d-restricted}
Even when each $|W_\zeta|$ is replaced by the size of its $D$-quotient
$|W_\zeta| / |W_\zeta \cap W_J|$ (restricting to the pairing parabolic coset
structure), the block sizes still disagree with the Rule B sizes across the
non-principal blocks. Enumerating $W_\zeta \subseteq S_6$ and the intersections with
$W_J = \langle s_1, s_3, s_5\rangle$ gives:
\begin{center}
\begin{tabular}{l|r|l|r|r}
$\zeta$ composition & $|W_\zeta|$ & $W_\zeta \cap W_J$ & $|W_\zeta \cap W_J|$ & $|W_\zeta|/|W_\zeta \cap W_J|$ \\
\hline
$(6)$ & $720$ & $W_J$ & $8$ & $90$ \\
$(5,1)$ positions $\{1..5\}$& $120$ & $\langle s_1, s_3\rangle$ & $4$ & $30$ \\
$(5,1)$ positions $\{2..6\}$& $120$ & $\langle s_3, s_5\rangle$ & $4$ & $30$ \\
$(4,2)$ positions $\{1..4\},\{5,6\}$& $48$ & $W_J$ & $8$ & $6$ \\
$(3,3)$ & $36$ & $\langle s_1, s_5\rangle$ & $4$ & $9$ \\
$(4,1,1)$ & $24$ & $\langle s_1, s_3\rangle$ & $4$ & $6$ \\
$(1,\ldots,1)$ & $1$ & $1$ & $1$ & $1$ \\
\end{tabular}
\end{center}
None of the $D$-quotient sizes equals $24, 18, 12$. So the Rule B block multiset
$\{30, 24, 18, 12, 6\}$ still does not match any assignment of $\zeta$'s to blocks.
\end{corollary}

\begin{proof}
Direct enumeration. (The relevant characters are listed by their ordered multiplicity
compositions; each orbit under $S_6$ gives a conjugate $W_\zeta$, and
$|W_\zeta|/|W_\zeta \cap W_J|$ depends on the embedding.)
\end{proof}

\section{Why the structures differ}

The mismatch is not incidental; it reflects the different categorical origins of the
two decompositions.

\begin{itemize}
\item \textbf{Rule B is value-indexed}: blocks are indexed by the value $w(n) = c$. This
is a \emph{left}-coset decomposition of $S_n$ by the point-stabilizer $S_{n-1}$ (of
value $n$), restricted to $D_n$. Equivalently by Lemma \ref{lem:rule-b-parabolic} it
is the \emph{right}-coset decomposition by the positional parabolic $W_k$. The two
descriptions coincide because $\Pi^{(k)} \in \mathcal{H}(W_k)$.

\item \textbf{Sandvik is character-indexed}: blocks are indexed by a torus character
$\zeta \in T$. The endoscopic subgroup $W_\zeta$ measures the symmetry of $\zeta$
under permutation of the torus coordinates. This is a finer data than a single value:
$\zeta$ specifies all multiplicities of distinguished torus eigenvalues, not one
position.
\end{itemize}

The one case where they align is $\zeta = (z, z, z, z, z, z')$ ($c=1$ frozen): the
endoscopic Levi $W_\zeta = S_5$ and the parabolic Levi $W_k = S_5$ coincide. Sandvik's
theorem then lifts the 30-dim principal block to a 120-dim categorical object by
re-adding the $|W_J \cap W_\zeta| = 4$-fold multiplicity killed by the $D$-restriction.
This lift exists, but it is not a \emph{proof} of the block identity — it is a
restatement.

\section{Consequence for $p_\lambda(q)$ palindromy}

The target of the Sandvik route was to derive the palindromic factor
$p_\lambda(q) = q^{\deg p_\lambda} p_\lambda(q^{-1})$ in the Rosetta Stone
eigenvalue formula from Sandvik's parity-sheaf self-Verdier-duality. Theorem
\ref{thm:mismatch} shows that Rule B is not a monodromic operation, so Sandvik's
categorical palindromy does not a-priori descend to $p_\lambda(q)$ via Rule B.

A direct alternative is already available:
\begin{observation}\label{obs:kl-route}
The parabolic Hecke module $\mathcal{H}_q \otimes_{\mathcal{H}_q(W_J)} \mathrm{sgn}_{W_J}$
carries a Kazhdan--Lusztig cellular structure; the image of the staircase element
$\Pi_q = \Pi^{(n-1)}$ on its $D_n$-basis is governed by this cellular structure.
Palindromy of per-irrep eigenvalue factors of $\Pi_q$ on $\mathbb{C}[D_n]$ follows
from the positivity of KL polynomials (already proved for this problem in
\cite{memory:eigenvalue-positivity-proved}) combined with the
$\mathbb{Z}/2$-symmetry of KL basis under the bar involution, which sends $T_w$ to
$q^{-\ell(w)} T_{w^{-1}}^{-1}$ and multiplies parabolic basis vectors by signs
$q^{\ell/2}$ compatible with self-conjugation.
\end{observation}

Observation \ref{obs:kl-route} is a suggestion, not a proof; its precise formalization
is the natural next step, replacing the Sandvik route.

\section{Computational verification}

All claims in this note were verified in SymPy at $n=6, k=4$ over $\mathbb{Q}(q)$:

\begin{enumerate}
\item $|D_6| = 90$, and the Rule B partition of $D_6$ by $w(6) \in \{1,\ldots,5\}$
gives block sizes $(30, 24, 18, 12, 6)$ (Proposition \ref{prop:block-sizes}).
\item The Rule B blocks equal the intersections of $D_6$ with the right
$W_4$-cosets of $S_6$ (Lemma \ref{lem:rule-b-parabolic}).
\item The principal block, reordered by the canonical bijection with $D_5$, equals
$M_R^{(4)}[D_5, D_5]$ entrywise as a $30 \times 30$ matrix of polynomials
(Proposition \ref{prop:principal}).
\item The Young-subgroup orders of $S_6$ are $\{1, 2, 4, 6, 8, 12, 24, 36, 48, 120, 720\}$,
ruling out the Rule B sizes $30, 18, 12$ as endoscopic orders (Lemma \ref{lem:young}
and Theorem \ref{thm:mismatch}).
\end{enumerate}

Scripts: \texttt{sandvik-block-compute.py} (block structure),
\texttt{sandvik-verify.py} (principal block identity and size comparison).
Both at \texttt{\~{}/projects/scratch/}.

\section{Status}

The Sandvik monodromic Hecke framework is ruled out as a direct categorical lift of
Block Decomposition Rule B. The remaining viable route for the palindromy of
$p_\lambda(q)$ is Observation \ref{obs:kl-route}: Kazhdan--Lusztig cellular structure
on the parabolic Hecke module, combined with the already-proved KL positivity for
the staircase matrix. Together with the Bi cluster negative result of
2026-04-21 \cite{memory:bi-cluster-n4-test}, this concentrates attention on the
KL / W-graph / Soergel-bimodule direction without the monodromic apparatus.

\begin{thebibliography}{9}

\bibitem{sandvik}
S.~Sandvik,
\emph{Monodromic Hecke 2-categories and endoscopic reduction},
arXiv:2604.15582, 2026.

\bibitem{memory:block-decomposition-rule-b}
Memory note \texttt{project\_block\_decomposition\_rule\_b.md}, 2026-04-18.

\bibitem{memory:eigenvalue-positivity-proved}
Memory note \texttt{project\_eigenvalue\_positivity\_proved.md}, 2026-04-17.

\bibitem{memory:bi-cluster-n4-test}
Memory note \texttt{project\_bi\_cluster\_n4\_test.md}, 2026-04-21.

\end{thebibliography}

\end{document}
